\documentclass[12pt, letterpaper]{article}

%% Build from the paper/ directory:
%%   cd paper && pdflatex main.tex
%% or use the root build system:
%%   make paper          (from repo root)
%%   ./build.sh paper    (POSIX)
%%   build.cmd paper     (Windows cmd)

% packages.tex
\usepackage{amsmath, amssymb, amsthm}
\usepackage{graphicx}
\usepackage{xcolor}
%% Tell the url package (loaded by hyperref) to allow line breaks
%% inside URLs at hyphens (in addition to the default / and .).
%% This prevents margin-drift on long paths like
%% dcl-paper-02-sm-derivation that would otherwise overflow.
\PassOptionsToPackage{hyphens}{url}
\usepackage[
    hidelinks,
    breaklinks=true,
    pdftitle={Geometry Forces Physics: A Lie-Algebra Derivation of the SM Gauge Group},
    pdfauthor={Jack Menendez},
    pdfsubject={A=1 Discrete Causal Lattice; SM gauge group; Lie-algebra automorphisms},
    pdfkeywords={Standard Model gauge group, automorphism algebra, bipartite causal lattice, SU(3), SU(2), U(1), SO(3,1), Lie algebra, symbolic verification}
]{hyperref}
\usepackage{booktabs}
\usepackage{array}
\usepackage{geometry}
\usepackage{adjustbox}
\usepackage{longtable}
\usepackage{setspace}
\usepackage{caption}
\usepackage{subcaption}
\usepackage{float}
\usepackage{tikz}
\usetikzlibrary{positioning, fit, calc, backgrounds, shapes.geometric, arrows.meta}
\onehalfspacing

\input{macros/commands}

%% Page geometry (also set here as a fallback in case the package path fails)
\geometry{top=1in, bottom=1in, left=1in, right=1in}

%% ---------------------------------------------------------------
%% TITLE PAGE
%% ---------------------------------------------------------------

\title{Geometry Forces Physics:\\
       A Lie-Algebra Derivation of the Standard Model Gauge Group\\
       from a Single Conservation Law%
       \thanks{Version 1.0.  Paper~II of the A=1 Discrete Causal
       Lattice series; direct continuation of
       \emph{Geometry First} (Paper~I,
       \href{https://doi.org/10.5281/zenodo.20078529}{DOI:
       10.5281/zenodo.20078529}).
       %% At release time, append the v1.0 Zenodo DOI:
       %% \href{https://doi.org/10.5281/zenodo.XXXXXXXX}{DOI: 10.5281/zenodo.XXXXXXXX}.
       }}
\author{Jack Menendez\\
        \small Eugene, OR\\
        \small ORCID: \href{https://orcid.org/0009-0003-1166-307X}{0009-0003-1166-307X}\\
        \small \href{mailto:jackdmenendez@gmail.com}{jackdmenendez@gmail.com}}
\date{\today}

\begin{document}

\maketitle

\begin{abstract}
    % abstract.tex
This is Paper~II of the A=1 Discrete Causal Lattice series. The
central thesis is to prove (or characterise the obstruction to)
Eq.~(137) of Paper~I: that the automorphism group of the bipartite
octahedral causal lattice $\mathcal{T}_\diamond^3$ under the unity
constraint $\mathcal{A}=1$ is the Standard Model gauge group plus
Lorentz, $SO(3,1) \times SU(3) \times SU(2) \times U(1)$, with real
dimension $18 = 6+8+3+1$.  We establish on the existing per-site
$\mathbb{C}^2 = (\psi_R, \psi_L)$ amplitude that $SO(3,1) \times U(1)$
acts as a dim-7 sub-product (verified symbolically), that the proposed
per-site $SU(2)$ generators on this carrier are not independent of
the Lorentz rotations, and that the RGB sublattice contributes only
$\mathbb{Z}_3 \subset SU(3)$.  We extend the framework to the per-site
amplitude $\mathbb{C}^{12} = \mathbb{C}^2 \otimes \mathbb{C}^2 \otimes
\mathbb{C}^3$ (chirality $\otimes$ weak isospin $\otimes$ colour),
verify that the four conjecture factors commute pairwise on this
carrier (dim 18), verify that a trivial tensor extension of the
bipartite tick rule preserves $\mathcal{A}=1$ and the global
$SU(2)_W \times SU(3)$ symmetries, and verify that the bipartite
Wilson plaquette $V_1, -V_2, -V_1, V_2$ is gauge-invariant.  The
remaining open subproblems --- exact equality versus containment in
$\mathrm{Aut}_\text{ext}$, SM-chirality coupling alignment, and the
$1/g^2$ prefactor for the non-abelian Wilson actions --- are stated
as the substance of the paper.  The paper closes with either a
derivation of the SM gauge group from a single conservation axiom on
a discrete substrate, or a precise statement of the obstruction; both
outcomes are concrete results.

\placeholder{(Trim and tighten as the paper develops; the audit table
in Section~\ref{sec:introduction} carries the canonical PASS/STUB
status for each sub-claim.)}

\end{abstract}

%% ---------------------------------------------------------------
%% FRONT-MATTER NOTICES (after abstract per author preference)
%% ---------------------------------------------------------------

\medskip
\noindent\textbf{Keywords.}  Standard Model gauge group; automorphism
algebra; bipartite causal lattice; SU(3); SU(2); U(1); SO(3,1);
direct product; Lie algebra; symbolic verification.

\medskip
\noindent\textbf{License.}  Paper text and figures: Creative Commons
Attribution 4.0 (CC BY 4.0).  Source code in the companion
repository: MIT License.

\medskip
\noindent\textbf{Data availability.}  Numerical experiments, data
files, and the LaTeX source for this paper are publicly available
at the GitHub repository
\nolinkurl{github.com/JackDMenendez/dcl-paper-02-sm-derivation}.
The Lie-algebra scaffolding scripts
\nolinkurl{automorphism_*.py} and \nolinkurl{tick_rule_*.py} were
introduced in Paper~I~\cite{menendez2026geometry} and are
reproduced here as the evidence base for the automorphism
conjecture; the seven additional sympy verification scripts that
close Phases~1-4 of the present work are new in this paper.
%% At release time, point to the tag and the Zenodo DOI:
%% The exact commit corresponding to this release is the v1.0 tag,
%% archived at
%% \href{https://doi.org/10.5281/zenodo.XXXXXXXX}%
%%      {\nolinkurl{doi:10.5281/zenodo.XXXXXXXX}}.

%% ---------------------------------------------------------------
%% RELEASE NOTES (only at release time)
%% ---------------------------------------------------------------
%% \bigskip
%% \noindent\textbf{Release notes for vX.Y (since vX.Y-1).}
%% \begin{itemize}
%% \item Item 1.
%% \item Item 2.
%% \end{itemize}

%% \emph{Status legend used throughout.}  \texttt{PASS} / \texttt{PART}
%% / \texttt{STUB} / \texttt{FAIL} -- see Table~\ref{tab:audit}.

\newpage
\tableofcontents
\newpage

%% ============================================================
%% PART I: <name>
%% ============================================================
\section{Introduction}
% introduction.tex (Paper II)
% Frames the central claim (Eq. (137) of Paper I) and includes the
% audit table (Table \ref{tab:audit}) so the reader sees the
% PASS/STUB decomposition before any dependent claim.

\label{sec:introduction}

Paper~I (\emph{Geometry First}) established the bipartite octahedral
causal lattice $\mathcal{T}_\diamond^3$ with the unity constraint
$\mathcal{A}=1$ as a discrete substrate from which the Dirac equation,
the Born rule, hydrogen quantization at four significant figures,
and the structural form of the induced gauge action are derived.
The paper closes (§15) with a precise central conjecture, repeated
here as the anchor of Paper~II:
%
\begin{equation}
\label{eq:automorphism_conjecture}
\mathrm{Aut}(\mathcal{T}_\diamond^3, \mathcal{A}=1)
\;=\; SO(3,1) \times SU(3) \times SU(2) \times U(1).
\end{equation}
%
Both sides are precise finite-dimensional Lie algebras with real
dimension $18 = 6 + 8 + 3 + 1$.  The match is more than dimensional:
the structure constants of the lattice's automorphism algebra must
reproduce those of the right-hand side factor by factor.  Either the
automorphism calculation closes the equality (and the Standard Model
gauge group is derived from a single conservation axiom on a discrete
substrate), or it does not (and Paper~II is a precise statement of
the obstruction).  Both outcomes are concrete results.

\subsection{Status of the conjecture entering Paper~II}
\label{subsec:status_at_entry}

Paper~I's v1.0 release supplies six computational scaffolding
scripts in \nolinkurl{src/utilities/}
(\nolinkurl{automorphism_*.py}, \nolinkurl{tick_rule_*.py}) that
establish the load-bearing
sub-claims of Eq.~\eqref{eq:automorphism_conjecture}.  On the
existing per-site amplitude $\mathbb{C}^2 = (\psi_R, \psi_L)$ the
direct product $SO(3,1) \times U(1)$ (dim 7) is verified, and the
proposed per-site $SU(2)$ generators on this carrier are confirmed
to coincide with the Lorentz rotation subgroup --- so
$SU(2)_W$ and $SU(3)$ require structural extensions of the per-site
amplitude.  On the extended amplitude
$\mathbb{C}^{12} = \mathbb{C}^2 \otimes \mathbb{C}^2 \otimes
\mathbb{C}^3$ the four conjecture factors commute pairwise (dim 18
verified), the bipartite tick rule extends consistently
($\mathcal{A}=1$, parity, and $SU(2)_W \times SU(3)$ all preserved),
and the bipartite Wilson plaquette $V_1, -V_2, -V_1, V_2$ is
gauge-invariant.  The audit table (Table~\ref{tab:audit}) records
each sub-claim with its supporting script.

%% The audit table is included immediately after the status section
%% so the reader meets it before any claim that depends on a status
%% entry.
% ============================================================
% audit_table.tex (Paper II)
%
% Maps Paper II's claims to the supporting derivations and to the
% sympy verification scripts in src/utilities/.
%
% The central claim is the conjecture stated in Paper~I as Eq.~(137):
%   Aut(T_diamond, A=1) = SO(3,1) x SU(3) x SU(2) x U(1)
% (real dim 18 = 6+8+3+1).  Paper II's task is to close this row by
% showing the four established sub-claims plus the three open
% sub-claims combine to the equality.
%
% Status semantics (paste into the abstract or front-matter as well):
%   PASS  -- derivation complete or verification confirms the claim.
%   PART  -- mechanism demonstrated, full quantitative match pending.
%   STUB  -- analytical result with no numerical confirmation yet,
%            or planned calculation not yet done.
%   FAIL  -- tested and disconfirmed.
%
% Use \texttt{} for code/script names and file paths inside cells.
% ============================================================

\label{tab:audit}

{\scriptsize
\setlength{\tabcolsep}{4pt}
\renewcommand{\arraystretch}{0.95}
\begin{longtable}{@{}p{3.2cm}p{3.6cm}p{3.4cm}p{3.8cm}p{1.0cm}@{}}
\caption{\small System audit: Paper~II claims mapped to supporting
derivations and sympy verification scripts.  Status semantics
(\texttt{PASS} / \texttt{PART} / \texttt{STUB} / \texttt{FAIL}) per
the front-matter.  The central claim is the bottom row;
its status is the conjunction of the rows above.}
\label{tab:audit_caption} \\
\hline
\textbf{Claim} & \textbf{Mechanism} & \textbf{Continuum / formal limit} & \textbf{Evidence} & \textbf{Status} \\
\hline
\endfirsthead

\multicolumn{5}{l}{\textit{(Table~\ref{tab:audit}, continued from previous page)}} \\
\hline
\textbf{Claim} & \textbf{Mechanism} & \textbf{Continuum / formal limit} & \textbf{Evidence} & \textbf{Status} \\
\hline
\endhead

\hline
\multicolumn{5}{r}{\textit{(continued on next page)}} \\
\endfoot

\hline
\endlastfoot

% ── Established factors ─────────────────────────────────────────────
Discrete spatial $\mathrm{Aut}(\Gamma, \mathbf{V})$ has order 48
    & Enumeration of basis-preserving linear maps on
      $\{\pm\mathbf{V}_1, \pm\mathbf{V}_2, \pm\mathbf{V}_3\}$
    & Cubic point group $B_3 \cong O_h$;
      12 elements orthogonal in $\mathbb{R}^3$
    & \nolinkurl{automorphism_discrete.py} (sympy)
    & \texttt{PASS} \\

RGB sublattice symmetry contributes only $\mathbb{Z}_3 \subset SU(3)$
    & Cyclic shift of the three RGB basis vectors generates a
      $\mathbb{Z}_3$ subgroup
    & Abelian; cannot generate non-abelian $SU(3)$ on its own
    & \nolinkurl{automorphism_rgb_su3.py} (sympy)
    & \texttt{PASS} \\

$SO(3,1) \times U(1)$ acts on the existing $\mathbb{C}^2 = (\psi_R, \psi_L)$
    & Lorentz from $O_h$-averaged Dirac (Paper~I §6); $U(1)$
      from the per-site phase oscillator
    & Real dim $6 + 1 = 7$, direct product on the existing
      per-site amplitude
    & \nolinkurl{automorphism_direct_product.py} (sympy);
      verifies $[\mathfrak{so}(3,1), \mathfrak{u}(1)] = 0$ and that
      the proposed per-site $SU(2)$ generators are the same matrices
      as the Lorentz rotation subgroup
    & \texttt{PASS} \\

Direct-product structure on extended $\mathbb{C}^{12}$
    & Per-site $\mathbb{C}^2 \otimes \mathbb{C}^2 \otimes \mathbb{C}^3$
      with the four factors acting on disjoint tensor slots
    & Real dim $6 + 8 + 3 + 1 = 18$;
      $\mathfrak{so}(3,1) \oplus \mathfrak{su}(3) \oplus
       \mathfrak{su}(2) \oplus \mathfrak{u}(1)$ commutes pairwise
    & \nolinkurl{automorphism_direct_product_extended.py} (sympy)
    & \texttt{PASS} \\

Tick-rule consistency on $\mathbb{C}^{12}$
    & Trivial tensor extension
      $T_\text{ext} = T_\text{chirality} \otimes I_2 \otimes I_3$
      of the bipartite tick rule
    & Preserves $\mathcal{A}=1$, bipartite parity, and the global
      $SU(2)_W \times SU(3)$ symmetries
    & \nolinkurl{tick_rule_extended_consistency.py} (sympy);
      four checks pass
    & \texttt{PASS} \\

Wilson plaquette gauge invariance
    & Bipartite Wilson plaquette
      $V_1, -V_2, -V_1, V_2$ (smallest non-trivial closed loop;
      no triangular plaquettes since $V_1 + V_2 + V_3 \neq 0$)
    & $\mathrm{Tr}(W) = 2\cos(a_1 + a_2 + a_3 + a_4)$
      gauge-invariant by cyclic property
    & \nolinkurl{tick_rule_gauge_invariance.py} (sympy);
      $U(1)$ explicit, $SU(2)$ sample, general non-abelian by
      $V V^\dagger = I$
    & \texttt{PASS} \\

SU(3) representation branch consistency
    & Only the $\mathbf{3} \oplus \mathbf{3}$ interpretation
      (Branch~A: $T_a = \lambda_a / 2$ on every site) commutes
      with the bipartite tick rule
      $T_\text{chir} \otimes I_2 \otimes I_3$; the
      $\mathbf{3} \oplus \bar{\mathbf{3}}$ interpretation
      (Branch~B: RGB in $\mathbf{3}$, CMY in $\bar{\mathbf{3}}$)
      obstructs on the 5 real Gell-Manns
      $\{\lambda_1, \lambda_3, \lambda_4, \lambda_6, \lambda_8\}$
    & Residual factors as $[T_\text{chir}, P_R] \otimes I_2 \otimes
      \mathrm{Re}(\lambda_a)$; forces the bipartite parity to be
      linear under the existing tick rule (antilinear $CP$
      candidates require modifying the tick rule itself)
    & \nolinkurl{su3_branch_consistency.py} (sympy);
      8/8 Branch~A commutators vanish, 3/8 Branch~B commute,
      5/8 Branch~B obstruct
    & \texttt{PASS} \\

Non-abelian $SU(3)$ from the $\mathbb{C}^3$ colour memory
    & The colour-memory tick rule applies $U_j$ on $\mathbb{C}^3$ on
      a tick along $\mathbf{V}_j$; the $S_3$ orbit $\{X_j = -i\log U_j\}$
      closes under Lie brackets to a subalgebra of $\mathfrak{su}(3)$.
      The natural ``rotate toward $|j\rangle$'' reading
      ($X_1 = \lambda_1 + \lambda_4$ for $j=1$) and any
      mixed Cartan-plus-real-off-diagonal candidate generate the
      full $\mathfrak{su}(3)$
    & Degenerate cases (purely diagonal Cartan only, or purely
      imaginary antisymmetric off-diagonals) yield strict
      subalgebras: $\mathfrak{h}$ (Cartan, dim 2) or
      $\mathfrak{su}(2)$ (dim 3) respectively.
      The discrete RGB $\mathbb{Z}_3$ of
      \nolinkurl{automorphism_rgb_su3.py} is recovered as
      the abelian-Cartan corner.  The $\mathfrak{su}(3)$ factor of
      Eq.~(137) is dynamically generated, not merely admitted
    & \nolinkurl{su3_generation_from_colour_memory.py} (sympy);
      8 candidate $X_1$'s surveyed, closure dimensions
      $\{0, 2, 2, 3, 3, 3, 8, 8\}$
    & \texttt{PASS} \\

\hline
% ── Open sub-claims ─────────────────────────────────────────────────
Exact equality vs containment in $\mathrm{Aut}_\text{ext}$
    & The discrete-Hermitian centralizer of the bipartite tick
      rule on $\mathbb{C}^{12}$ (= operators $X$ commuting with
      $\sigma_x \otimes I_2 \otimes I_3$) has dimension 71 in
      $\mathfrak{su}(12)$, equivalently $\mathfrak{u}(6) \oplus
      \mathfrak{u}(6) - 1$.  Of the 14 Hermitian generators of
      Eq.~(137), exactly 12 ($J_1$, 3 $SU(2)_W$, 8 $SU(3)_c$)
      lie in the centralizer; the remaining 2 ($J_2, J_3$ of the
      Lorentz rotations) are continuum-emergent
    & 59 extra generators commute with the tick rule but are
      \emph{factor-mixing} operators -- couplings like
      $I_2 \otimes \sigma_a \otimes \lambda_b$ (isospin-colour mixing,
      24 generators), $\sigma_x \otimes \sigma_a \otimes I_3$
      (3), $\sigma_x \otimes I_2 \otimes \lambda_a$ (8), and
      $\sigma_x \otimes \sigma_a \otimes \lambda_b$ (24).
      These violate the SM's factor-product gauge invariance.
      Equality $=$ in Eq.~(137) holds iff the factor-product
      structure is imposed as an additional constraint
    & \nolinkurl{automorphism_centralizer_extended.py} (sympy);
      143 $\mathfrak{su}(12)$ generators enumerated, exactly 71
      commute, 12 are conjecture-aligned and 59 are factor-mixings
    & \texttt{PART} \\

SM-chirality coupling alignment
    & Under Branch~A SU(3) and the existing tick rule, the unique
      linear bipartite parity is $P = \sigma_x \otimes I_2 \otimes I_3$
      (by Schur on the irreducible $\mathbf{3}$).
      $P$ anticommutes with $\gamma_5 = \sigma_z \otimes I_2 \otimes I_3$
      and swaps the SM's $\psi_L \leftrightarrow \psi_R$ (spatial parity,
      not chirality projection);
      $P$ commutes with the proposed $SU(2)_W$ generators (vector-like,
      not the SM's chiral coupling).
      Bipartite $\mathbb{Z}_2$ and SM chirality $\mathbb{Z}_2$ are
      orthogonal Bloch involutions on the chirality $\mathbb{C}^2$
    & SM left-handed doublets / right-handed singlets are
      empirically required; framework provides instead a vector-like
      $SU(2)$ coupling and a spatial-parity $\mathbb{Z}_2$ that is
      orthogonal to (not equal to) the SM's chirality $\mathbb{Z}_2$.
      The follow-on question (whether a modified antilinear tick
      rule admits SM chirality / $CP$) is settled negatively below
    & \nolinkurl{chirality_parity_alignment.py} (sympy);
      5 candidates tested, unique survivor $M = I_3$ (and its
      projective equivalent $-I_3$); all 8 structural identities
      verified
    & \texttt{PART} \\

Modified tick rule for Branch~B SU(3) compatibility
    & Ansatz $T_\text{ext}^B = is \cdot I_{12} + c \cdot (\sigma_x
      \otimes I_2 \otimes C) \circ K$ with $K$ global complex
      conjugation: requires the colour matrix $C$ to anticommute with
      every Gell-Mann generator $\lambda_a$
    & No non-zero $C$ exists: anticommutation with $\lambda_3$
      restricts $C$ to entries $C_{12}, C_{21}, C_{33}$;
      anticommutation with $\lambda_8$ then forces those to zero.
      Representation-theoretically: $\mathbf{3}$ and $\bar{\mathbf{3}}$
      of SU(3) are inequivalent irreps; the unique antilinear
      intertwiner is global colour conjugation, which does not
      satisfy the chirality-dependent commutativity Branch~B
      demands.  Lattice cannot incorporate SM-style $CP$ as a
      discrete symmetry of any natural tick modification
    & \nolinkurl{tick_rule_cp_modified.py} (sympy);
      9 candidates tested, all fail; structural proof via
      $\lambda_3, \lambda_8$ anticommutation;
      isospin entanglement does not relax the colour constraint
    & \texttt{FAIL} \\

Explicit $1/g^2$ prefactor for $SU(2)_W$, $SU(3)$ Wilson actions
    & Non-abelian generalisation of Paper~I's bipartite-plaquette
      calculation: $1 - \tfrac{1}{N}\mathrm{Re}\,\mathrm{Tr}\,W_{ab} =
      \tfrac{a^4 T_F}{2N}(V_a^i V_b^j F^c_{ij})^2$ with $T_F = 1/2$.
      The Q-tensor (eigenvalues $\{4, 4, 16\}$) and the
      $O_h$-averaged Maxwell-style density are inherited from
      Paper~I unchanged
    & Spectator-factor counting on the per-site
      $\mathbb{C}^2 \otimes \mathbb{C}^2 \otimes \mathbb{C}^3$
      gives $N_f^{SU(2)} = 6$, $N_f^{SU(3)} = 4$; sharp lattice-
      scale prediction $g_3^2/g_2^2 = 3/2$ independent of the
      universal one-loop prefactor $c$ Paper~I left open.
      Universal $c$ remains open (the explicit
      $-\mathrm{Tr}\ln D_\text{lat}[U]$ calculation is a
      paper-length follow-on for both papers)
    & \nolinkurl{induced_gauge_action_nonabelian.py} (sympy);
      Q-tensor verified, spectator counts derived, ratio
      $g_1^2 : g_2^2 : g_3^2 = 1 : 4 : 6$ at lattice scale
    & \texttt{PART} \\

\hline
% ── Central claim ───────────────────────────────────────────────────
\textbf{Eq.~(137) of Paper~I:}
$\mathrm{Aut}(\mathcal{T}_\diamond^3, \mathcal{A}=1) = SO(3,1) \times SU(3) \times SU(2) \times U(1)$
    & Conjunction of the four established factors above plus the
      three open sub-claims; structure constants must match factor
      by factor, not just dimensions
    & Standard Model gauge group plus Lorentz, derived from a
      single conservation axiom on a discrete substrate
    & Resolves favourably iff the three open sub-claims resolve
      favourably; characterises the obstruction otherwise
    & \texttt{STUB} \\

\hline
\end{longtable}
}


\subsection{Open subproblems}
\label{subsec:open_subproblems}

What Paper~II takes up:

\begin{enumerate}
\item \textbf{Exact equality vs containment} in
$\mathrm{Aut}_\text{ext}$.  The verified statement is
$\supseteq$; the conjecture asserts $=$.  Whether the extended
structure admits additional symmetries beyond the four-factor
product is a finite Lie-algebra enumeration question constrained by
commutation with the established factors.

\item \textbf{SM-chirality coupling alignment.}  The Standard Model's
$SU(2)_W$ couples only to left-handed fermion doublets; right-handed
fermions are $SU(2)_W$ singlets.  Whether the framework's bipartite
RGB/CMY parity is the SM's chirality projector, and whether the
proposed $SU(2)_W$ on the per-site weak-isospin $\mathbb{C}^2$
couples asymmetrically to $\psi_R$ versus $\psi_L$ as required by
the SM, is the load-bearing structural-alignment problem.

\item \textbf{Explicit $1/g^2$ prefactor} for the $SU(2)_W$ and
$SU(3)$ Wilson actions on the bipartite octahedral lattice.  The
existing $U(1)$ induced-gauge-action calculation (Paper~I
Appendix~B) supplies the template; the non-abelian generalisation
goes through with matrix-valued link variables.
\end{enumerate}

\subsection{Scope and outcomes}
\label{subsec:scope}

Paper~II's scope is precisely the three subproblems above.  All
other open questions noted in Paper~I --- the Farey-sequence mass
spectrum (Paper~I follow-on \#3), the Veneziano connection
(\#3), the proton-interior structure (\#16 dependency), plasma /
gradient ionization (\#1), and the like --- are out of scope and
deferred to subsequent papers in the series.

The paper has two viable terminal states.  If the three
subproblems resolve in the same direction, the conjecture closes and
the Standard Model gauge group is derived; the paper publishes that
derivation.  If any subproblem fails, the paper publishes a precise
statement of the obstruction and the framework's verified scope
(\emph{relativistic Hilbert-space + gravity + Born-rule structures}
plus $SO(3,1) \times U(1)$ on the existing per-site $\mathbb{C}^2$,
per Paper~I) is restated unchanged.


\section{Established Factors}
% established_factors.tex (Paper II)
% Structural prose for the PASS rows of audit_table.tex.  Each
% subsection takes one PASS row and explains what the script's output
% means structurally -- i.e. WHY the verified statement is the right
% statement, not just THAT it is verified.  The closing subsection
% (\autoref{subsec:established_necessity_vs_sufficiency}) collects the
% rows into the statement "established containment $\supseteq$"
% and names the gap to "equality $=$" that the open subproblems
% (\autoref{subsec:open_subproblems}) close.

\label{sec:established_factors}

The audit table (Table~\ref{tab:audit}) records the PASS rows
establishing sub-claims of Eq.~\eqref{eq:automorphism_conjecture},
together with one structural-consistency row (SU(3) representation
branch) that narrows the interpretation of the colour factor.  Each
is a sympy-verified statement on the bipartite octahedral lattice
$\mathcal{T}_\diamond^3$.  This section explains what these rows
mean structurally and how they combine into the containment
$\mathrm{Aut}(\mathcal{T}_\diamond^3, \mathcal{A}=1) \supseteq SO(3,1)
\times SU(3) \times SU(2) \times U(1)$ on the extended per-site
amplitude $\mathbb{C}^{12} = \mathbb{C}^2 \otimes \mathbb{C}^2
\otimes \mathbb{C}^3$.  The closing subsection
(\autoref{subsec:established_necessity_vs_sufficiency}) states the
gap between the verified containment and the conjectured equality;
that gap is the substance of Paper~II's three open subproblems
(\autoref{subsec:open_subproblems}).

\subsection{Discrete spatial automorphisms: $|\mathrm{Aut}_\text{discrete}(\Gamma, \mathbf{V})| = 48$}
\label{subsec:discrete_spatial}

The basis $\mathbf{V} = \{\pm\mathbf{V}_1, \pm\mathbf{V}_2,
\pm\mathbf{V}_3\}$ of the bipartite octahedral lattice has six
elements.  The linear maps $A \in GL(3,\mathbb{R})$ permuting
$\mathbf{V}$ as a set are parametrised by a permutation $\sigma \in
S_3$ of $(\mathbf{V}_1, \mathbf{V}_2, \mathbf{V}_3)$ together with
independent signs $s \in \{\pm 1\}^3$ on each basis vector.  This
gives $|S_3| \cdot |\mathbb{Z}_2^3| = 6 \cdot 8 = 48$ distinct
candidates, all of which preserve $\mathbf{V}$ as a set.  The script
\nolinkurl{automorphism_discrete.py} enumerates the 48 candidates,
verifies each preserves $\mathbf{V}$, checks closure under matrix
multiplication, and confirms that exactly 12 of the 48 elements
satisfy $A^\top A = I$ in the standard inner product on
$\mathbb{R}^3$.

The resulting group is the hyperoctahedral group
$B_3 \cong \mathbb{Z}_2 \wr S_3$, abstractly the cubic point group
$O_h$.  The realisation on $\mathbb{R}^3$ in the $\mathbf{V}$-basis
is not, however, the standard cubic-symmetry orthogonal
representation: only the 12-element orthogonal subgroup is metric-
preserving, while the remaining 36 elements are linear shears that
permute the basis labels at the cost of distorting angles.

Structurally, the full 48-element group is the relevant discrete
spatial symmetry for the framework: the bipartite tick rule treats
the RGB and CMY labels symmetrically and does not single out the
standard metric.  The 12-element orthogonal subgroup is the
metric-preserving part --- the genuine crystallographic point
symmetry of $\mathcal{T}_\diamond^3$ in the standard
$\mathbb{R}^3$ embedding --- and the standard metric is restored
in the continuum limit by the $O_h$-averaging argument of
Paper~I~§6 that produces $SO(3,1)$.  No piece of the 48 is
discarded; the non-orthogonal piece is the discrete record of the
fact that the lattice does not commit to a Euclidean metric at the
microscale.

\subsection{The RGB sublattice contributes only $\mathbb{Z}_3 \subset SU(3)$}
\label{subsec:rgb_z3_obstruction}

The discrete RGB symmetry $S_3$ of $(\mathbf{V}_1, \mathbf{V}_2,
\mathbf{V}_3)$ embeds in $U(3)$ as the six permutation matrices on
the kinetic-step weight space $(\alpha_1, \alpha_2, \alpha_3) \in
\mathbb{C}^3$.  Within $U(3)$, the determinant separates the six
permutations: the three even permutations have $\det = +1$ and lie
in $SU(3)$ as the alternating subgroup $A_3 \cong \mathbb{Z}_3$;
the three transpositions have $\det = -1$ and lie in $U(3) \setminus
SU(3)$.  The script \nolinkurl{automorphism_rgb_su3.py} verifies the
embedding and the determinant split explicitly.

This is the load-bearing structural observation that forces the
per-site $\mathbb{C}^3$ extension.  The conjecture's $SU(3)$ factor
has real dimension 8 and is non-abelian; $\mathbb{Z}_3$ is finite
and abelian; closure under multiplication does not extend a finite
abelian discrete group to a continuous non-abelian Lie group, and
the commutator $[A, B] = AB - BA$ on the three $\mathbb{Z}_3$
matrices either vanishes or lands back in $\mathbb{Z}_3$, so the
Lie algebra generated is itself abelian.  An $\mathfrak{su}(3)$
factor cannot be a Lie-theoretic continuation of the framework's
existing discrete colour symmetry.  It must come from an additional
per-site internal index that the existing $\mathbb{C}^2 = (\psi_R,
\psi_L)$ amplitude does not carry.

The natural candidate is a per-site $\mathbb{C}^3$ ``colour-memory''
amplitude $(c_1, c_2, c_3)$ where component $c_j$ records the
amplitude that the wavefunction's most recent RGB tick was along
$\mathbf{V}_j$.  The eight Gell-Mann generators
$\{\lambda_a/2\}_{a=1\ldots 8}$ acting on this $\mathbb{C}^3$
supply $\mathfrak{su}(3)$ by definition; the framework's
contribution is identifying this $\mathbb{C}^3$ as the
\emph{minimal} per-site internal extension the conjecture's
$SU(3)$ factor needs.  The existing $\mathbb{Z}_3$ remains as the
discrete subgroup of $SU(3)$ that permutes the colour basis ---
not refuted but recovered as a finite subgroup of the larger
continuous symmetry.

\subsection{$SO(3,1) \times U(1)$ on the existing per-site $\mathbb{C}^2$: dim 7}
\label{subsec:lorentz_u1_on_C2}

Four candidate Lie-algebra factors of the conjecture admit
generators on the existing per-site chirality amplitude
$\mathbb{C}^2 = (\psi_R, \psi_L)$: the Lorentz rotations
$J_a = \tfrac{1}{2}\sigma_a$, the Lorentz boosts $K_a =
\tfrac{i}{2}\sigma_a$, a proposed per-site $SU(2)$ with generators
$T_a = \tfrac{1}{2}\sigma_a$, and the per-site $U(1)$ phase
generator $Q = i I_2$.  The script
\nolinkurl{automorphism_direct_product.py} verifies that
$\mathfrak{so}(3,1)$ closes correctly on $(J_a, K_a)$ (the standard
Weyl-spinor relations $[J_a, J_b] = i\epsilon_{abc}J_c$,
$[K_a, K_b] = -i\epsilon_{abc}J_c$, $[J_a, K_b] = i\epsilon_{abc}
K_c$) and that $Q$ commutes with every generator.  So
$\mathfrak{so}(3,1) \oplus \mathfrak{u}(1)$ acts as a direct sum
on $\mathbb{C}^2$, with total real dimension $6 + 1 = 7$.

The script also verifies the literal matrix identity $T_a = J_a$
for every $a \in \{1, 2, 3\}$.  The proposed per-site $SU(2)$
generators on the chirality $\mathbb{C}^2$ are not just isomorphic
to the Lorentz rotation generators --- they are the same matrices.
As subgroups of $GL(2, \mathbb{C})$, the proposed per-site $SU(2)$
on this carrier is the Lorentz rotation subgroup of $SO(3,1)$.  A
direct-product decomposition $SO(3,1) \times SU(2)$ on the
existing $\mathbb{C}^2$ would double-count the rotations.

This is the second load-bearing structural observation, and it
parallels the $\mathbb{Z}_3 \subset SU(3)$ finding: the
conjecture's $SU(2)$ factor, like its $SU(3)$ factor, cannot be a
direct continuation of the framework's existing structure.  An
independent per-site $\mathbb{C}^2$ internal index --- naturally
identified with weak isospin --- is required for $SU(2)_W$ to
contribute a separate factor.  Without that extension the
framework's verified per-site carrier supports only $SO(3,1)
\times U(1)$ at dim 7, half the conjecture's dim 18.

\subsection{Direct-product structure on the extended $\mathbb{C}^{12}$: dim 18}
\label{subsec:direct_product_C12}

The previous two subsections identified the minimal per-site
extension required by the conjecture: the per-site amplitude
\begin{equation}
\label{eq:c12_amplitude}
\psi[\text{chir}, \text{iso}, \text{col}] \in
\mathbb{C}^2_\text{chir} \otimes \mathbb{C}^2_\text{iso} \otimes
\mathbb{C}^3_\text{col} \;\cong\; \mathbb{C}^{12},
\end{equation}
with the four conjecture factors acting on disjoint tensor slots:
\begin{align*}
SO(3,1):\;\;& J_a \otimes I_2 \otimes I_3,\quad
              K_a \otimes I_2 \otimes I_3,\\
SU(2)_W:\;\; & I_2 \otimes \tfrac{1}{2}\sigma_a \otimes I_3,\\
SU(3):\;\;   & I_2 \otimes I_2 \otimes \tfrac{1}{2}\lambda_a,\\
U(1):\;\;    & i\,I_{12}.
\end{align*}
The script \nolinkurl{automorphism_direct_product_extended.py}
verifies the within-factor commutation relations
($\mathfrak{so}(3,1)$, $\mathfrak{su}(2)$, $\mathfrak{su}(3)$ all
close correctly) and the pairwise commutation
$[\mathfrak{so}(3,1), \mathfrak{su}(2)] = [\mathfrak{so}(3,1),
\mathfrak{su}(3)] = [\mathfrak{su}(2), \mathfrak{su}(3)] = 0$ by
the tensor identity $[A \otimes I, I \otimes B] = 0$.  The
$\mathfrak{u}(1)$ generator $i I_{12}$ is central and commutes with
everything.

This establishes the direct-sum Lie algebra
\begin{equation}
\label{eq:direct_sum_18}
\mathfrak{so}(3,1) \oplus \mathfrak{su}(3) \oplus
\mathfrak{su}(2) \oplus \mathfrak{u}(1)
\;\subseteq\; \mathfrak{aut}(\mathcal{T}_\diamond^3,
\mathcal{A}=1)
\end{equation}
on the extended $\mathbb{C}^{12}$ carrier, with real dimension
$6 + 8 + 3 + 1 = 18$.  Both the dimension and the direct-sum
structure of the right-hand side of
Eq.~\eqref{eq:automorphism_conjecture} are now realised on the
extended amplitude.  What remains for the conjecture's
\emph{equality} is the question of whether the inclusion in
\eqref{eq:direct_sum_18} is proper --- whether the extended
$\mathrm{Aut}_\text{ext}$ admits any generator beyond the
direct-sum 18-dimensional Lie algebra constructed here.  That is
the open subproblem of
\autoref{subsec:open_subproblems}, item (i).

\subsection{Tick-rule consistency on the extended $\mathbb{C}^{12}$}
\label{subsec:tick_rule_consistency}

The bipartite chirality tick operator
\[
T_\text{chir} = \begin{pmatrix}
  i\sin(\delta\phi/2) & \cos(\delta\phi/2) \\
  \cos(\delta\phi/2)  & i\sin(\delta\phi/2)
\end{pmatrix}
\]
on $\mathbb{C}^2 = (\psi_R, \psi_L)$ is unitary and preserves
$\mathcal{A} = 1$ per Paper~I.  The natural extension to the
$\mathbb{C}^{12}$ carrier of
Eq.~\eqref{eq:c12_amplitude} is the trivial tensor extension
\begin{equation}
\label{eq:trivial_tensor_extension}
T_\text{ext} \;=\; T_\text{chir} \otimes I_2 \otimes I_3,
\end{equation}
acting as the existing tick rule on chirality and as the identity
on isospin and colour.  The script
\nolinkurl{tick_rule_extended_consistency.py} verifies four
structural properties:

\begin{enumerate}
\item \textbf{Unitarity}: $T_\text{ext}^\dagger T_\text{ext} = I_{12}$,
so $\mathcal{A}=1$ on $\mathbb{C}^{12}$ follows from $\mathcal{A}=1$
on the chirality $\mathbb{C}^2$ by linearity.
\item \textbf{Global $SU(2)_W$ commutativity}: $[T_\text{ext}, I_2
\otimes \tfrac{1}{2}\sigma_a \otimes I_3] = 0$ for $a = 1, 2, 3$.
\item \textbf{Global $SU(3)$ commutativity}: $[T_\text{ext}, I_2
\otimes I_2 \otimes \tfrac{1}{2}\lambda_a] = 0$ for $a = 1, \ldots, 8$.
\item \textbf{Bipartite parity}: with parity operator $P_\text{ext}
= \sigma_x \otimes I_2 \otimes I_3$, the identity $P_\text{ext}
T_\text{ext} P_\text{ext} = T_\text{ext}$ holds --- the chirality
parity that exchanges $\psi_R \leftrightarrow \psi_L$ is a
symmetry of the extended tick rule, and the isospin and colour
factors are parity-singlets in the trivial extension.
\end{enumerate}

The structural reading: the trivial tensor extension does no harm
to the framework's existing invariants.  $\mathcal{A}=1$, bipartite
parity, and the global $SU(2)_W \times SU(3)$ symmetry all
survive.  What this does \emph{not} establish is that the trivial
extension has physical content beyond six independent copies of the
existing chirality dynamics.  Promoting the global $SU(2)_W \times
SU(3)$ symmetry to a local gauge symmetry --- and identifying the
chirality-asymmetric coupling that the Standard Model demands --- is
the substance of the open subproblems in
\autoref{subsec:open_subproblems}, items (ii)--(iii).  The
established content here is precisely that the global symmetry is
\emph{available}; the gauge dynamics that turns it into physics is
the open work.

\subsection{Wilson plaquette gauge invariance on the bipartite lattice}
\label{subsec:wilson_plaquette}

Gauging the global $SU(2)_W \times SU(3)$ of
\autoref{subsec:tick_rule_consistency} introduces link variables
$U_i(x) \in G = U(1) \times SU(2)_W \times SU(3)$ on each
basis-vector edge, transforming as
\[
\psi(x) \to V(x)\psi(x),\qquad
U_i(x) \to V(x)\, U_i(x)\, V^\dagger(x + \mathbf{V}_i)
\]
under a local gauge transformation $V(x) \in G$.  The link-dressed
matter bilinear
\[
\psi^\dagger(x)\, U_i(x)\, \psi(x + \mathbf{V}_i)
\]
is gauge-invariant by direct substitution; the script
\nolinkurl{tick_rule_gauge_invariance.py} verifies this explicitly
for $U(1)$ and by sample check for $SU(2)$ link variables, with the
general non-abelian case (including $SU(3)$) following from
$V V^\dagger = I$.

The smallest closed gauge-invariant loop is structurally distinctive
on the bipartite octahedral lattice.  The three RGB basis vectors
satisfy $\mathbf{V}_1 + \mathbf{V}_2 + \mathbf{V}_3 = (1, 1, -1)
\neq 0$, so there are no triangular plaquettes; the smallest
closed loop respecting the RGB/CMY alternation is the four-link
``bipartite plaquette''
\[
x \to x + \mathbf{V}_1 \to x + \mathbf{V}_1 - \mathbf{V}_2 \to
x - \mathbf{V}_2 \to x,
\]
traversing the basis vectors in the cyclic order $\mathbf{V}_1,
-\mathbf{V}_2, -\mathbf{V}_1, \mathbf{V}_2$.  The plaquette variable
$W_p = U_1(x)\, U_{-2}(x + \mathbf{V}_1)\, U_{-1}(x + \mathbf{V}_1
- \mathbf{V}_2)\, U_2(x - \mathbf{V}_2)$ has gauge-invariant trace
by the cyclic property of $\mathrm{Tr}$ --- the gauge
transformations $V(x), V(x + \mathbf{V}_1), V(x + \mathbf{V}_1 -
\mathbf{V}_2), V(x - \mathbf{V}_2)$ at the four vertices cycle and
cancel.  This holds for any $SU(N)$ gauge group, in particular for
the $SU(2)_W$ and $SU(3)$ of the conjecture.

The Wilson action template
\begin{equation}
\label{eq:wilson_action_template}
S_W \;=\; \frac{1}{g^2} \sum_p \left(1 - \frac{1}{N}
\mathrm{Re}\,\mathrm{Tr}\, W_p\right)
\end{equation}
is therefore gauge-invariant on the bipartite octahedral lattice
by inspection.  The framework's existing $U(1)$ induced-action
calculation (Paper~I, Appendix~B) is the worked example of how the
bare coupling $1/g^2$ is extracted at the lattice scale; the
template \eqref{eq:wilson_action_template} extends to $SU(2)_W$ and
$SU(3)$ with matrix-valued link variables and no other structural
change.  Computing the explicit $1/g^2(a)$ prefactor for the
non-abelian cases is the substance of the open subproblem of
\autoref{subsec:open_subproblems}, item (iii).

\subsection{Necessity vs sufficiency: established containment $\supseteq$, conjectured equality $=$}
\label{subsec:established_necessity_vs_sufficiency}

The six PASS sub-claims of
Eq.~\eqref{eq:automorphism_conjecture} combine as follows; the
seventh PASS row (SU(3) representation branch consistency) narrows
the colour interpretation that the combination assumes (see
\nolinkurl{notes/su3_branch_consistency.md}).
\autoref{subsec:discrete_spatial} establishes the discrete spatial
backbone $O_h$.  \autoref{subsec:rgb_z3_obstruction} and
\autoref{subsec:lorentz_u1_on_C2} together force the per-site
internal extension $\mathbb{C}^2_\text{iso} \otimes
\mathbb{C}^3_\text{col}$: $SU(2)_W$ cannot share the chirality
$\mathbb{C}^2$ with $SO(3,1)$ rotations, and $SU(3)$ cannot be a
Lie-theoretic continuation of the discrete RGB symmetry.
\autoref{subsec:direct_product_C12} verifies that the extended
$\mathbb{C}^{12}$ carrier supports the four conjecture factors as
a direct-sum 18-dimensional Lie algebra with the structure
constants of $\mathfrak{so}(3,1) \oplus \mathfrak{su}(3) \oplus
\mathfrak{su}(2) \oplus \mathfrak{u}(1)$.
\autoref{subsec:tick_rule_consistency} verifies that the bipartite
tick rule extends consistently to $\mathbb{C}^{12}$ as a global
symmetry of the dynamics.  \autoref{subsec:wilson_plaquette}
verifies that the global symmetry can be promoted to a local gauge
symmetry with a gauge-invariant Wilson plaquette on the smallest
non-trivial closed loop the bipartite lattice admits; the explicit
$1/g^2$ prefactor that turns the gauge-invariant plaquette into a
full Wilson action remains open (item (iii) below).

The conjunction is the inclusion of
Eq.~\eqref{eq:direct_sum_18}: the conjecture's right-hand side is
contained in the lattice's automorphism algebra,
\begin{equation}
\label{eq:established_containment}
\mathrm{Aut}(\mathcal{T}_\diamond^3, \mathcal{A}=1)
\;\supseteq\;
SO(3,1) \times SU(3) \times SU(2) \times U(1),
\end{equation}
with both sides finite-dimensional Lie groups of real dimension
$\geq 18$ on the extended $\mathbb{C}^{12}$ carrier.

The containment in \eqref{eq:established_containment} is necessary
but not sufficient for the conjectured equality
\eqref{eq:automorphism_conjecture}.  Three structural questions
remain open, in the order they enter the equality:
\begin{enumerate}
\item \emph{Equality vs proper containment in
$\mathrm{Aut}_\text{ext}$}: whether the extended automorphism
algebra admits any generator beyond the 18 already constructed.
This is a finite Lie-algebra enumeration question constrained by
commutation with the established factors.

\item \emph{Chirality coupling alignment} (partially resolved,
characterisation outcome): whether the framework's bipartite parity
(the RGB/CMY $\mathbb{Z}_2$) is the SM's chirality projector, and
whether the proposed $SU(2)_W$ on the isospin $\mathbb{C}^2$
couples to $\psi_R$ and $\psi_L$ asymmetrically as the Standard
Model demands.  Under Branch~A SU(3) and the existing tick rule,
the unique linear bipartite parity $P = \sigma_x \otimes I_2 \otimes
I_3$ anticommutes with $\gamma_5 = \sigma_z \otimes I_2 \otimes I_3$
and swaps $\psi_L \leftrightarrow \psi_R$ --- it IS spatial parity,
not chirality projection.  The proposed $SU(2)_W$ commutes with $P$
(vector-like, not the SM's chiral coupling).  The lattice's
bipartite $\mathbb{Z}_2$ and the SM's chirality $\mathbb{Z}_2$ are
distinct orthogonal Bloch involutions on the chirality
$\mathbb{C}^2$.  The trivial tensor extension
\eqref{eq:trivial_tensor_extension} therefore predicts a precise
\emph{characterisation} of the parity / chirality relationship
rather than an alignment or a vacuous obstruction.  The follow-on
question --- whether a modified antilinear tick rule of the form
$T_\text{ext}^B = is \cdot I_{12} + c \cdot (\sigma_x \otimes I_2
\otimes C) \circ K$ can admit the Branch~B SU(3) interpretation and
deliver SM-style $CP$ --- is settled negatively: no $3 \times 3$
matrix $C$ anticommutes with every Gell-Mann, by direct argument on
$\lambda_3$ and $\lambda_8$, equivalently because $\mathbf{3}$ and
$\bar{\mathbf{3}}$ are inequivalent irreps of SU(3) and the unique
antilinear intertwiner between them (global colour conjugation)
fails to satisfy the chirality-dependent commutativity Branch~B
demands.  The lattice's bipartite parity is therefore intrinsically
\emph{linear} (spatial parity, no charge conjugation), and the
framework cannot incorporate SM-style $CP$ as a discrete symmetry
of any natural tick modification.

\item \emph{Explicit $1/g^2$ prefactors for the non-abelian Wilson
actions}: generalising Paper~I's $U(1)$ induced-gauge-action
calculation to $SU(2)_W$ and $SU(3)$ with matrix-valued links, and
checking the resulting bare couplings against the measured
Standard-Model values at the lattice scale.
\end{enumerate}
These three questions are the substance of
\autoref{subsec:open_subproblems}.  Their joint resolution decides
between the conjecture's two terminal states: equality (Paper~II
publishes the derivation of the SM gauge group) or precise
characterisation of the obstruction (Paper~II publishes the
framework's verified scope plus the structural reason the equality
fails).  Both are concrete results; the established containment
\eqref{eq:established_containment} is what makes either of them
publishable.


%% Add more \section{...} \input{sections/...} blocks as needed.
%% Group with comment dividers (e.g., "PART II: DYNAMICS") to keep
%% main.tex skimmable as the paper grows.

\section{Conclusion}
% conclusion.tex
%
% Phase 5 synthesis (v1.0).  The conclusion is structured as:
%   - subsec:results_summary         -- what each phase established
%   - subsec:terminal_characterisation -- the framework's terminal prediction
%   - subsec:quantitative_predictions  -- the testable numbers
%   - subsec:methodological           -- substrate-first, debt to measurement
%   - subsec:open_directions          -- follow-on papers and open questions

\label{sec:conclusion}

\subsection{What Phase~1 through Phase~4 established}
\label{subsec:results_summary}

Eq.~\eqref{eq:automorphism_conjecture} of Paper~I is settled to
\texttt{PART} (characterisation outcome).  The four phases of the
work plan (\nolinkurl{notes/work_plan.md}) close as follows.

\paragraph{Phase~1 -- Chirality alignment (\texttt{PART}).}  Under
Branch~A SU(3) and the existing bipartite tick rule, the unique
linear bipartite parity is $P = \sigma_x \otimes I_2 \otimes I_3$.
It anticommutes with $\gamma_5 = \sigma_z \otimes I_2 \otimes I_3$,
swaps the SM's $\psi_L \leftrightarrow \psi_R$, and commutes with
every proposed $SU(2)_W$ generator.  The lattice's bipartite
$\mathbb{Z}_2$ and the SM's chirality $\mathbb{Z}_2$ are therefore
orthogonal Bloch involutions on the chirality $\mathbb{C}^2$:
spatial parity, not chirality projection.  Symbolically verified
by \nolinkurl{src/utilities/chirality_parity_alignment.py}.

\paragraph{Phase~1.5 -- Modified tick rule for SM $CP$ (\texttt{FAIL}).}
The natural antilinear modification
$T_\text{ext}^B = is\,I_{12} + c\,(\sigma_x \otimes I_2 \otimes C) \circ K$
that would make Branch~B SU(3) ($\mathbf{3} \oplus \bar{\mathbf{3}}$) a
global symmetry requires the colour matrix $C$ to anticommute with
every Gell-Mann generator.  No non-zero $C$ does: anticommutation
with $\lambda_3$ leaves only $C_{12}, C_{21}, C_{33}$ unconstrained,
and anticommutation with $\lambda_8$ kills those.  The lattice
cannot incorporate SM-style $CP$ as a discrete symmetry of any
natural tick modification.  Symbolically verified by
\nolinkurl{src/utilities/tick_rule_cp_modified.py}; the
representation-theoretic origin is that $\mathbf{3}$ and
$\bar{\mathbf{3}}$ are inequivalent irreps of SU(3).

\paragraph{Phase~2 -- Non-abelian $SU(3)$ from $\mathbb{C}^3$ memory (\texttt{PASS}).}
The natural rotate-toward-$|j\rangle$ colour-memory tick rule
($X_1 = \lambda_1 + \lambda_4$ for $j = 1$, with $S_3$ conjugates)
closes under Lie-bracket closure to the full $\mathfrak{su}(3)$.
Restricted candidates (diagonal-only, single-off-diagonal) yield
strict subalgebras (Cartan $\mathfrak{h}$ or one $\mathfrak{su}(2)$
subalgebra).  The framework's $\mathfrak{su}(3)$ factor is
\emph{dynamically generated} by the lattice for any natural
real-symmetric colour-memory tick rule; the discrete RGB
$\mathbb{Z}_3$ of \nolinkurl{automorphism_rgb_su3.py} is recovered as
the abelian-Cartan corner of this larger structure.  Symbolically
verified by \nolinkurl{src/utilities/su3_generation_from_colour_memory.py}.

\paragraph{Phase~3 -- Containment vs equality on $\mathrm{Aut}_\text{ext}$ (\texttt{PART}).}
The discrete-Hermitian centralizer of the bipartite tick rule on
$\mathbb{C}^{12}$ is
\begin{equation}
\label{eq:su6_su6_u1}
\mathfrak{Aut}_{\text{centralizer}} \;=\; \mathfrak{su}(6)_+ \oplus \mathfrak{su}(6)_- \oplus \mathfrak{u}(1),
\qquad \dim = 35 + 35 + 1 = 71,
\end{equation}
with the two $\mathfrak{su}(6)$ factors acting on the $\pm 1$
chirality eigenspaces and $J_1$ as the central $\mathfrak{u}(1)$.
The SM gauge subalgebra $\mathfrak{su}(3) \oplus \mathfrak{su}(2)
\oplus \langle J_1 \rangle$ (dim 12 in $\mathfrak{su}(12)$) embeds
diagonally via the $\mathbf{6} = (\mathbf{3}, \mathbf{2})$
branching; the 59 ``extras'' decompose under $SU(3) \times SU(2)$
as
\begin{equation}
\label{eq:extras_decomp}
\text{59 extras} \;=\; 48 \times (\mathbf{8}, \mathbf{3})_{\text{leptoquark}}
\;+\; 11 \times \text{chirality-}\sigma_x\text{ shadow SM}.
\end{equation}
Bracket-structure verification
(\nolinkurl{aut_centralizer_extras_commutators.py}):
$[\text{extras}, \text{SM}] \subseteq \text{extras}$ (0 of 708
brackets land in SM); $[\text{extras}, \text{extras}]$ has 226 of
1711 brackets in SM, so the extras are an SM-module but not a Lie
ideal --- exactly the algebraic shape of a left-right symmetric
broken gauge symmetry pattern.  Symbolically verified by
\nolinkurl{src/utilities/automorphism_centralizer_extended.py}.

\paragraph{Phase~4 -- Wilson $1/g^2$ for $SU(2)_W$, $SU(3)$ (\texttt{PART}).}
Paper~I's bipartite-plaquette Q-tensor (eigenvalues
$\{4, 4, 16\}$) is inherited unchanged for non-abelian gauge
groups; the SU(N) trace identity introduces only a global
$T_F = 1/2$ factor.  Spectator-factor counting on the per-site
$\mathbb{C}^{12}$ gives the framework's quantitative gauge-coupling
prediction:
\begin{equation}
\label{eq:lattice_scale_ratios}
\frac{1}{g_1^2} : \frac{1}{g_2^2} : \frac{1}{g_3^2} \;=\; 12 : 3 : 2 \quad \text{at the lattice scale } 1/a,
\end{equation}
equivalently $g_1^2 : g_2^2 : g_3^2 = 1 : 4 : 6$, and the sharp
ratio
\begin{equation}
\label{eq:g3_over_g2}
g_3^2 / g_2^2 = 3/2 \quad \text{at the lattice scale,}
\end{equation}
independent of the universal one-loop prefactor $c$ (which
Paper~I and Paper~II both leave open as the explicit
$-\mathrm{Tr}\,\ln D_\text{lat}[U]$ calculation).  Symbolically
verified by \nolinkurl{src/utilities/induced_gauge_action_nonabelian.py}.

\subsection{The framework's terminal characterisation}
\label{subsec:terminal_characterisation}

The conjunction of Phase~1 through Phase~4 settles the conjecture's
central row to \texttt{PART} with a precise structural reading.
Eq.~\eqref{eq:automorphism_conjecture} holds as containment
\begin{equation}
\label{eq:containment_again}
\mathrm{Aut}(\mathcal{T}_\diamond^3, \mathcal{A}=1)
\;\supseteq\;
SO(3,1) \times SU(3) \times SU(2) \times U(1),
\end{equation}
where the right-hand side is the \emph{factor-product-irreducible
subalgebra} of the lattice's discrete-Hermitian automorphism
centralizer of Eq.~\eqref{eq:su6_su6_u1}.  Equality $=$ does
\emph{not} hold at the discrete level: the centralizer is 71-dim,
not 18-dim; the SM gauge group is exactly the projection that
respects the tensor-factor decomposition
$\mathbb{C}^{12} = \mathbb{C}^2_\text{chir} \otimes
\mathbb{C}^2_\text{iso} \otimes \mathbb{C}^3_\text{col}$.

The framework's terminal prediction:

\begin{itemize}
\item The lattice's natural symmetry at the discrete tick scale
is the left-right symmetric algebra
$\mathfrak{su}(6)_+ \oplus \mathfrak{su}(6)_- \oplus \mathfrak{u}(1)$
of Eq.~\eqref{eq:su6_su6_u1}, with the bipartite parity acting as
the $L \leftrightarrow R$ exchange.  This is the algebraic shape
of a left-right symmetric grand unification, derived from the
$\mathcal{A}=1$ constraint on the bipartite octahedral lattice
rather than postulated.

\item The SM gauge group $SU(3) \times SU(2) \times U(1)$ is the
factor-product subalgebra of this larger structure, recovered
under the (interpretive) restriction that gauge invariance respects
the tensor decomposition.  The 59 ``extras''
(Eq.~\eqref{eq:extras_decomp}) are coset generators that
transform as $(\mathbf{8}, \mathbf{3})$ leptoquark-flavoured plus
chirality-shadow SM under the SM adjoint action.

\item The lattice's bipartite parity is intrinsically linear
(\emph{spatial} parity), and the proposed $SU(2)_W$ couples
\emph{vector-like} to matter, not chirally.  The SM's chirality
and $CP$ violation are not derivable from the framework at the
discrete level --- they are external choices that the SM's
Lagrangian formalism imposes by separating kinetic and mass terms.
The framework's mass mechanism (mass-as-clock-density,
\nolinkurl{notes/mass_chirality_coupling.md}) is structurally
incompatible with chiral gauge coupling.

\item Gravity is a scalar clock-density field on the flat
bipartite lattice (Paper~I~§7), not a curved spacetime metric.
Spacetime torsion --- the antisymmetric part of the affine
connection in Einstein-Cartan theory --- is therefore not a
separate ingredient the framework requires; its conventional
motivation (coupling Dirac spin to curved gravity) is moot
because the framework provides the Dirac structure intrinsically
through the bipartite tick rule's chirality oscillation, without
ever invoking a spin connection (\nolinkurl{notes/no_spacetime_torsion.md}).
\end{itemize}

\subsection{Quantitative predictions}
\label{subsec:quantitative_predictions}

Three quantitative predictions emerge from the Phase~1-4 work:

\paragraph{Gauge-coupling ratio at the lattice scale.}
Eq.~\eqref{eq:g3_over_g2}: $g_3^2 / g_2^2 = 3/2$ at $1/a$.  The SM
measured ratio at $M_Z$ is $g_3^2/g_2^2 \approx 3.3$; the
factor-of-$\sim 2$ shift is consistent with RG running over many
energy decades (asymptotic freedom of $g_3$, slow running of
$g_2$).  The framework does not yet fix the lattice scale $1/a$ from
first principles, so the prediction is a ratio at an unknown UV
scale; a follow-on paper would extract $1/a$ either from the
explicit one-loop prefactor calculation or from another framework
observable.

\paragraph{Anisotropic photon dispersion along $(1, 1, -1)$ (Paper~I); structural extension to non-abelian.}
Paper~I~§6 / \nolinkurl{induced_gauge_action.tex} established
gauge-sector birefringence aligned with the optical axis
$\mathbf{V}_1 + \mathbf{V}_2 + \mathbf{V}_3 = (1, 1, -1)$ for the
abelian photon, derived from the bipartite-plaquette Q-tensor
with eigenvalues $\{4, 4, 16\}$.  Phase~4's non-abelian
generalisation inherits the same Q-tensor unchanged, so the
non-abelian gauge bosons ($W^\pm$, $Z$, gluons) inherit the same
direction-dependent kinetic structure as a structural corollary
of the bipartite geometry.  We flag this as a structural
consequence, not as an independent quantitative prediction
verified in this paper.

\paragraph{Anisotropic decoherence in the colour-weak sector (hypothesis).}
A structural hypothesis (\nolinkurl{notes/extras_as_a1_accounting.md},
\texttt{DRAFT}): the 59 extras of Eq.~\eqref{eq:extras_decomp} are
the algebraic bookkeeping channels through which $\mathcal{A}=1$
tracks probability during what an SM observer identifies as
decoherence.  Decoherence, in this reading, is the algebraic
projection from the full 71-dim centralizer onto the SM-gauge-
invariant 12-dim subalgebra; the ``lost coherence'' resides in
the 59 extras directions.  The
$(\mathbf{8}, \mathbf{3})$ representation content of the
leptoquark-flavoured extras predicts colour-octet weak-triplet
anisotropy in decoherence rates --- a specific testable departure
from continuum decoherence theory.  Verification requires the
quantitative follow-on calculation sketched in
\nolinkurl{notes/extras_as_a1_accounting.md}; this paper does not
execute it.

\subsection{Methodological reflections}
\label{subsec:methodological}

The framework's contribution is sharper than ``a derivation of the
SM gauge group from a discrete substrate.''  Paper~II identifies
\emph{precisely which} features of the Standard Model are geometric
consequences of $\mathcal{A}=1$ on $\mathcal{T}_\diamond^3$ and
\emph{which are not}.

\paragraph{Geometric:} the Lie algebra structure
$\mathfrak{so}(3,1) \oplus \mathfrak{su}(3) \oplus \mathfrak{su}(2)
\oplus \mathfrak{u}(1)$ (dim 18) embedded in a larger left-right
symmetric algebra (dim 71); the bipartite plaquette gauge
invariance and the Q-tensor with eigenvalues $\{4, 4, 16\}$; the
discrete spatial $O_h$ point group (order 48); the colour
$\mathfrak{su}(3)$ generated dynamically from the
$\mathbb{C}^3$ memory.

\paragraph{Not geometric (external SM choices):} chirality (the SM
selects $\sigma_z$ as the chirality axis where the lattice has
$\sigma_x$ as bipartite parity --- orthogonal Bloch involutions);
$CP$ violation (no antilinear modification of the lattice tick
rule admits a Branch~B $CP$); the Higgs mechanism separating
kinetic and mass terms (the lattice's mass \emph{is} the kinetic
mechanism, via the chirality-mixing $\cos(\delta\phi/2)$
amplitude); the factor-product gauge invariance (the lattice
admits 59 factor-mixing operators that the SM forbids by
construction).

\paragraph{The substrate-first approach.}  The framework begins
from a single conservation axiom on a discrete substrate and
works out what physics follows.  This is different from
constructing a model around observed phenomena and fitting
parameters; the framework's predictions are derived from the
lattice geometry without post-hoc parameter fitting.  Paper~I
(\href{https://doi.org/10.5281/zenodo.20078529}{DOI:
10.5281/zenodo.20078529}; see its audit table for the PASS rows)
demonstrated this by deriving the Dirac equation, the Born rule,
hydrogen quantization at four significant figures, and the
structural form of the induced gauge action.  Paper~II continues
the programme, identifying the Lie-algebra-level features that
the substrate produces and the Lagrangian-level features it does
not.

\paragraph{The debt to measurement.}  That the framework's
derivations land where measurement has spent a century pinning
the answers down is not a coincidence --- it is the prerequisite.
A century of accurate, well-tested physics is what makes
substrate-first derivation feasible in the first place: there
must be something firm to land on.  The framework's contribution
is not to replace conventional physics but to identify which of
its features are geometric consequences of a single conservation
law and which are not.  Both readings honour what the testing
community established; the framework reads the established answers
as emergent rather than postulated, but it has no answers to
offer about features the established answers have not already
pinned down (\nolinkurl{notes/debt_to_measurement.md}).

\paragraph{The audit-table convention.}  The audit-table discipline
inherited from Paper~I (Section~\ref{subsec:audit_convention}) is
what makes documentation of emergence credible rather than
aspirational.  Every claim of the paper traces to an artifact:
a sympy-verified statement in \nolinkurl{src/utilities/}, a
derivation in the text, or a measured value cited from the
literature.  Drift between drafts is mechanically detectable;
negative results stay visible (Phase~1.5's \texttt{FAIL} row on
SM-style $CP$ is part of the paper's contribution, not an
embarrassment to be elided); the canonical record is the table,
not the prose.  The convention is a small mechanical change with a
substantial accountability shift, and we hope it is useful to
other first-principles theoretical proposals.

\subsection{Open directions and follow-on papers}
\label{subsec:open_directions}

Several follow-on calculations are scoped but not executed in
this paper.

\paragraph{The universal one-loop prefactor $c$.}  Paper~I left
the explicit $-\mathrm{Tr}\,\ln D_\text{lat}[U]$ calculation as a
``paper-length follow-on''; Paper~II inherits this open item.
Closing it would fix the absolute scale of $1/g^2(1/a)$, enabling a
quantitative RG-running check against the measured
$g_i^2(M_Z)$ values.

\paragraph{The lattice scale $1/a$.}  The framework's lattice-scale
predictions assume a UV cutoff that is not yet fixed from first
principles.  Possible inputs: the Planck scale, the hydrogen
quantization scale, or the framework's induced gravitational
coupling.  Fixing $1/a$ converts the ratio prediction
$g_3^2/g_2^2 = 3/2$ into a quantitative test against measured
couplings at any energy scale via RG running.

\paragraph{The decoherence calculation.}  The structural
hypothesis captured in \nolinkurl{extras_as_a1_accounting.md}
converts the 59 extras into a predictive structure for
anisotropic decoherence rates.  A follow-on script
\nolinkurl{decoherence_as_projection.py} (sketched in
the note) would compute the rate of probability flow from
SM-projected to extras-projected wavefunctions as a function of
the tick rule's $\delta\phi$ and the initial coherence,
converting the qualitative hypothesis into a quantitative
prediction.

\paragraph{The symmetry-breaking mechanism.}  The 71-dim
centralizer at the lattice scale reduces to the SM's 12-dim
gauge group at observable energies.  By analogy with standard GUT
breaking, the 59 extras would correspond to massive gauge bosons
that decouple at lower energies.  Identifying the framework's
breaking mechanism --- whether it is Higgs-analogous, decoherence-
driven via the algebraic projection of
\nolinkurl{notes/extras_as_a1_accounting.md}, or geometric in
some other sense --- is the natural next paper in the series.

\paragraph{Follow-on papers in the catalogue.}  The downstream
papers indexed in
\nolinkurl{external/dcl/notes/follow_on_implications.md} that depend
on Paper~II's colour machinery include item~\#16's proton-
interior study, item~\#3's Farey-sequence mass spectrum, and
items \#13 (Operation Algebra of the DCL) and \#14 (Balanced
Equations and Birefringent Channels) on which the framework's
notational and algebraic infrastructure rests.  Paper~II's
results enter those papers as imported infrastructure: the SU(3)
representation content of colour-memory amplitudes, the
left-right symmetric extension, and the quantitative
gauge-coupling ratios at the lattice scale.

\subsection{Closing}
\label{subsec:closing}

Eq.~\eqref{eq:automorphism_conjecture} of Paper~I asked whether the Standard Model gauge group
plus Lorentz emerges from a single conservation axiom on the
bipartite octahedral causal lattice.  The answer, after four
phases of work: yes as containment, with a precise structural
characterisation of the gap to equality.  The framework realises
a left-right symmetric algebra $\mathfrak{su}(6) \oplus
\mathfrak{su}(6) \oplus \mathfrak{u}(1)$ at the lattice scale,
with the SM gauge group as its factor-product projection.  The
framework's quantitative prediction $g_3^2/g_2^2 = 3/2$ at the
lattice scale, the bipartite-plaquette birefringence aligned with
the $(1, 1, -1)$ optical axis, and the anisotropic decoherence
hypothesis are the testable consequences this work hands to
experiment.

What the framework derives is the Lie-algebra structure of the SM
gauge group from substrate geometry; what it does not derive is
the SM's chirality, $CP$ violation, or Higgs mechanism.  Both
readings are concrete contributions: the geometric features are
recovered without parameter fitting, and the non-geometric
features are identified as external choices in the SM's
Lagrangian formulation that the substrate does not impose.  The
discrete causal lattice with $\mathcal{A}=1$ does not reach the
full Standard Model, but it reaches the maximal structure a
substrate-first derivation can produce, and it identifies the
external inputs the SM requires beyond the substrate.


\section*{Acknowledgements}
% acknowledgements.tex

The author thanks Claude (Anthropic, Claude Opus 4.7 with 1M-token
context) for collaborative work throughout Paper~II's Phase~1--5
development: drafting the symbolic verification scripts in
\nolinkurl{src/utilities/}, writing substantial portions of the
technical prose and working notes (\nolinkurl{notes/}), engaging
with structural questions about the framework, and maintaining the
audit-table claim-status discipline across drafts.  The Phase~3
identification of the discrete-Hermitian centralizer as
$\mathfrak{su}(6) \oplus \mathfrak{su}(6) \oplus \mathfrak{u}(1)$
(sharpening an earlier ``71-dimensional centralizer'' statement
into the explicit left-right symmetric structure) was prompted by
a structural review from Google's Gemini model, shared by the
author in collaborative discussion.  All claims, decisions, and
the final form of the paper remain the author's editorial and
scientific responsibility.

The author also thanks the broader physics community whose century
of careful measurement is the foundation on which substrate-first
derivation is feasible at all
(\nolinkurl{notes/debt_to_measurement.md}): the predictions of
Paper~I (the Dirac equation, the Born rule, hydrogen quantization
at four significant figures, the structural form of the induced
gauge action) and Paper~II ($g_3^2/g_2^2 = 3/2$ at the lattice
scale, the bipartite-plaquette birefringence inherited unchanged
to the non-abelian sector) land against measured values that
spectroscopy and high-energy physics have pinned down through
testing accumulated over generations.  The framework owes its
verifiability to that empirical groundwork.

No external funding to acknowledge.


%% ============================================================
%% APPENDICES
%% ============================================================
\appendix

\section{Code and Data}
% code_and_data.tex -- Appendix A
\label{sec:code_and_data}

This appendix describes the companion repository
\nolinkurl{github.com/JackDMenendez/dcl-paper-02-sm-derivation}
(v1.0 tag), with the aim that any reader can audit or reproduce
any claim in the paper from the repository contents.  The
audit-table convention of Section~\ref{subsec:audit_convention}
applies: every claim has a row in Table~\ref{tab:audit} with a
status tag and a pointer to its supporting artifact (a script, a
derivation, or a measured value cited from the literature).

\subsection{Repository layout}
\label{subsec:repo_layout}

The repository is organised as follows.

\begin{itemize}
\item \nolinkurl{paper/} --- LaTeX source for this paper.
  \begin{itemize}
  \item \nolinkurl{paper/main.tex} --- the document root.
  \item \nolinkurl{paper/sections/} --- one file per section,
    including \nolinkurl{audit_table.tex} (the canonical
    record of every PASS / PART / STUB / FAIL claim status).
  \item \nolinkurl{paper/macros/} --- shared LaTeX packages and
    commands.
  \item \nolinkurl{paper/paper-bib/references.bib} --- the
    bibliography source.
  \end{itemize}

\item \nolinkurl{src/utilities/} --- sympy verification scripts.
  Thirteen scripts in total: six inherited from Paper~I (the
  \nolinkurl{automorphism_*.py} and \nolinkurl{tick_rule_*.py}
  files) and seven new to Paper~II that close Phases~1--4 of
  \nolinkurl{notes/work_plan.md}.  Each script is self-contained
  and runs from a Python virtual environment with \texttt{sympy}
  (see Appendix~\ref{sec:reproducibility} for setup).  The output
  of each run is the operational PASS/PART/FAIL for the
  corresponding audit-table row; running every script is the
  end-to-end verification of Paper~II's claims.

\item \nolinkurl{notes/} --- working theoretical notes that the
  paper either cites (via the audit table's evidence column) or
  builds on.  At v1.0 the directory contains:
  \begin{itemize}
  \item \nolinkurl{work_plan.md} --- the phased plan for Phases~0
    through 5, with resolution stanzas for each closed phase.
  \item Per-phase findings and structural interpretations
    (one note per closed Phase or sub-phase):
    \begin{itemize}
    \item \nolinkurl{su3_branch_consistency.md}
    \item \nolinkurl{chirality_alignment.md}
    \item \nolinkurl{cp_modification_obstruction.md}
    \item \nolinkurl{su3_generation_from_colour_memory.md}
    \item \nolinkurl{aut_centralizer_enumeration.md}
    \item \nolinkurl{aut_centralizer_extras_commutators.md}
    \item \nolinkurl{induced_gauge_action_nonabelian.md}
    \end{itemize}
  \item Structural-synthesis and framing notes:
    \begin{itemize}
    \item \nolinkurl{mass_chirality_coupling.md}
    \item \nolinkurl{debt_to_measurement.md}
    \item \nolinkurl{no_spacetime_torsion.md}
    \end{itemize}
  \item \nolinkurl{extras_as_a1_accounting.md} (\texttt{DRAFT})
    --- the hypothesis that the 59 extras of the Phase~3
    centralizer are the algebraic bookkeeping channels for
    $\mathcal{A}=1$ during what an SM observer identifies as
    decoherence.  Sketches the follow-on script
    \nolinkurl{decoherence_as_projection.py}, deferred from v1.0.
  \item \nolinkurl{lie_algebra_proof_sketch.md} --- a pointer to
    Paper~I's working proof sketch
    (\nolinkurl{external/dcl/notes/lie_algebra_automorphism_proof_sketch.md},
    accessible via the Windows directory junction described in
    \nolinkurl{CLAUDE.md}).
  \end{itemize}

\item \nolinkurl{release_notes/} --- per-version change logs and
  GitHub Release bodies (see also
  Appendix~\ref{subsec:repro_release_flow}).

\item \nolinkurl{CLAUDE.md} --- working brief: project status,
  what NOT to change, cross-references to Paper~I and the parallel
  formalisation effort in physics-research, notes index.  This
  file is the entry point for anyone (human or AI) returning to
  the project.

\item \nolinkurl{CITATION.cff} --- machine-readable citation
  metadata.

\item \nolinkurl{external/} --- Windows directory junctions to
  Paper~I (\nolinkurl{external/dcl}) and the physics-research
  formalisation effort (\nolinkurl{external/research}), gitignored
  per \nolinkurl{.gitignore}.  On a fresh clone these need to be
  recreated locally to follow cross-references to those
  repositories; the seeded prose explicitly cites Paper~I's
  Zenodo DOI~\cite{menendez2026geometry} so that a reviewer
  without those local checkouts can still navigate the published
  artifact.
\end{itemize}

\subsection{The audit table as canonical record}
\label{subsec:audit_canonical}

The convention introduced in Section~\ref{subsec:audit_convention}
is the load-bearing one for code/data audit: prose claims defer
to Table~\ref{tab:audit}, every row has a status tag, every
status tag has a pointer to a verifiable artifact, and the table
itself is part of the released artifact (frozen at v1.0).  A
reader who doubts any specific claim in the paper can navigate to
the relevant audit-table row, follow the evidence-column pointer
to the supporting script, and re-run the script for independent
verification.  Negative results (\texttt{FAIL} rows --- Phase~1.5
in particular) are part of the contribution and are recorded with
the same convention.

The current v1.0 audit-table state is:

\begin{itemize}
\item Established factors (\texttt{PASS}, $8$ rows): discrete
  spatial $\mathrm{Aut}(\Gamma, \mathbf{V}) = O_h$ of order 48;
  RGB sublattice contributes only $\mathbb{Z}_3 \subset SU(3)$;
  $SO(3,1) \times U(1)$ acts on the existing $\mathbb{C}^2$
  (dim 7); direct-product structure on extended $\mathbb{C}^{12}$
  (dim 18); tick-rule consistency on $\mathbb{C}^{12}$; bipartite
  Wilson plaquette gauge invariance; SU(3) representation branch
  consistency; non-abelian $SU(3)$ from $\mathbb{C}^3$ memory.
\item Partial / characterisation (\texttt{PART}, $3$ rows):
  SM-chirality coupling alignment, exact equality vs containment
  in $\mathrm{Aut}_\text{ext}$, explicit $1/g^2$ prefactor for
  the non-abelian Wilson actions.
\item Disconfirmed (\texttt{FAIL}, $1$ row): modified tick rule
  for Branch~B SU(3) (no antilinear modification admits SM-style
  $CP$).
\item Central claim (\texttt{PART}): Eq.~(137) of Paper~I --
  containment $\supseteq$ holds; equality $=$ does not; the
  characterisation outcome is the framework's terminal
  prediction.
\end{itemize}

\subsection{Provenance and license}
\label{subsec:provenance}

Paper text and figures are released under the Creative Commons
Attribution 4.0 (CC BY 4.0) licence.  Source code in the
\nolinkurl{src/utilities/} directory is released under the MIT
licence.  The six inherited scripts (\nolinkurl{automorphism_*.py}
and \nolinkurl{tick_rule_*.py}) are verbatim copies of the v1.0
versions in Paper~I~\cite{menendez2026geometry}; edits to those
scripts in Paper~II are additive (extensions of the verification,
not refactorings of the existing calculation) so that Paper~I's
printed output remains the canonical reference for the original
six rows.

The complete commit history is available at the repository's
\nolinkurl{main} branch; the v1.0 release tag pins the exact
commit corresponding to this paper's published artifact.  Issues
and pull requests are accepted at the GitHub repository.


\section{Reproducibility Quickstart}
% reproducibility.tex -- Appendix B
\label{sec:reproducibility}

This appendix is a compact reproducibility quickstart.  Goal: a
reader who has just cloned the repository at the v1.0 tag can
reproduce every audit-table row of Paper~II and rebuild this
PDF from a fresh checkout in under ten minutes.

\subsection{Environment}
\label{subsec:repro_environment}

\begin{itemize}
\item \textbf{Operating system.}  Verified on Windows~11 (MSYS2
  UCRT64) and any POSIX system with the equivalent toolchain.
  No platform-specific code.
\item \textbf{Python.}  Python~3.10 or later.  The framework's
  symbolic verification scripts use only \texttt{sympy} (no
  numerical-only dependencies); the dependency manifest is
  \nolinkurl{virtual-env-requirements.txt} in the repository
  root.
\item \textbf{LaTeX.}  Any modern TeX distribution (TeX Live~2023+
  or MiKTeX) with \texttt{pdflatex} and standard packages
  (\texttt{amsmath}, \texttt{hyperref}, \texttt{longtable},
  \texttt{caption}, \texttt{adjustbox}; see
  \nolinkurl{paper/macros/packages.tex}).
\item \textbf{GNU Make}~$\geq$ 4.3 (the stock Windows GNU~Make port
  is 3.81 and too old; use the \texttt{chocolatey/make} build or
  MSYS2's \texttt{make}).
\end{itemize}

\subsection{From fresh clone to verified paper}
\label{subsec:repro_steps}

\begin{verbatim}
git clone https://github.com/JackDMenendez/dcl-paper-02-sm-derivation
cd dcl-paper-02-sm-derivation
git checkout v1.0

python -m venv .venv-ucrt64
source .venv-ucrt64/bin/activate    # POSIX / MSYS2
# .venv-ucrt64\Scripts\activate     # Windows native cmd
pip install -r virtual-env-requirements.txt

make paper                          # rebuilds build/Paper.pdf
\end{verbatim}

The \texttt{make paper} target runs three passes of
\texttt{pdflatex} (plus a \texttt{bibtex} pass for the
bibliography), producing
\nolinkurl{build/Paper.pdf} byte-equivalent to the released
artifact archived in
\nolinkurl{.stage/Paper_v1.0.pdf}.

\subsection{Verifying every audit-table row}
\label{subsec:repro_audit}

The end-to-end verification of Paper~II's PASS / PART / FAIL
status rows is running the thirteen scripts in
\nolinkurl{src/utilities/}.  Each script is self-contained and
runs as a Python module:

{\scriptsize
\begin{verbatim}
python -m src.utilities.automorphism_discrete
python -m src.utilities.automorphism_rgb_su3
python -m src.utilities.automorphism_direct_product
python -m src.utilities.automorphism_direct_product_extended
python -m src.utilities.tick_rule_extended_consistency
python -m src.utilities.tick_rule_gauge_invariance
python -m src.utilities.su3_branch_consistency
python -m src.utilities.chirality_parity_alignment
python -m src.utilities.tick_rule_cp_modified
python -m src.utilities.su3_generation_from_colour_memory
python -m src.utilities.automorphism_centralizer_extended
python -m src.utilities.aut_centralizer_extras_commutators
python -m src.utilities.induced_gauge_action_nonabelian
\end{verbatim}
}

All but \nolinkurl{aut_centralizer_extras_commutators.py} complete
in seconds on a consumer-grade laptop.  The bracket-structure
script (1711 commutators of $12\times 12$ sympy matrices) takes
several minutes; this is the most expensive computation in the
paper and is the load-bearing verification for the structural
finding that the 71-dim discrete centralizer is the standard
subalgebra-with-complement of a simple semisimple Lie algebra
(see~\autoref{subsec:terminal_characterisation}).

\subsection{Verifying the central numerical claim}
\label{subsec:repro_central}

The framework's first quantitative gauge-coupling prediction
($g_3^2 / g_2^2 = 3/2$ at the lattice scale, Eq.~\eqref{eq:g3_over_g2})
is derived in \nolinkurl{src/utilities/induced_gauge_action_nonabelian.py}.
The relevant lines extract the prediction from the Q-tensor and
the spectator-factor counting:

{\scriptsize
\begin{verbatim}
import sympy as sp
from sympy import Rational
# Spectator-factor multiplicities on the per-site C^{12}:
N_f_SU2 = 2 * 3   # chirality x colour spectators = 6
N_f_SU3 = 2 * 2   # chirality x isospin spectators = 4
T_F = Rational(1, 2)  # fundamental of SU(N)
inv_g2_sq = N_f_SU2 * T_F   # = 3
inv_g3_sq = N_f_SU3 * T_F   # = 2
print(inv_g2_sq / inv_g3_sq)   # 3/2 (= g_3^2 / g_2^2 at lattice scale)
\end{verbatim}
}

The full script also verifies Paper~I's Q-tensor with eigenvalues
$\{4, 4, 16\}$, the $T_F$ trace normalisation, and the abelian
($U(1)$) case where all 12 components of the per-site amplitude
carry unit charge, giving $1/g_1^2 : 1/g_2^2 : 1/g_3^2 = 12 : 3 : 2$
at the lattice scale.

\subsection{Where to look next}
\label{subsec:repro_pointers}

\begin{itemize}
\item For the per-claim verification trail: navigate from any
  audit-table row in Table~\ref{tab:audit} to the script named in
  its evidence column.  Run that script.
\item For structural-interpretation context: the corresponding
  note in \nolinkurl{notes/} contains the working derivation, the
  representation-theoretic decomposition, and pointers to adjacent
  notes (the upstream-flow tags at the bottom of each note name
  the related findings).
\item For the framework's foundational substrate:
  Paper~I~\cite{menendez2026geometry}.  Section~5--6 of Paper~I
  derive the Dirac structure from the bipartite tick rule;
  Section~7 derives gravity-as-clock-density; Appendix~B
  establishes the bipartite-plaquette Q-tensor and the
  structural form of the induced gauge action.  Paper~II's
  Phase~4 calculation extends Paper~I's Appendix~B unchanged to
  the non-abelian case.
\item For follow-on questions and open items: the catalogue at
  \nolinkurl{external/dcl/notes/follow_on_implications.md}
  (Paper~I repository) enumerates the planned downstream papers;
  Paper~II is item~\#16 of that catalogue.
\end{itemize}

\subsection{Release flow (for reference)}
\label{subsec:repro_release_flow}

The release procedure for v1.0 (and subsequent versions) is
documented in \nolinkurl{release_notes/README.md} of the
repository.  Briefly: (i)~final commits land on \texttt{main};
(ii)~the change log
(\nolinkurl{release_notes/v1.0.md}) and the GitHub Release body
(\nolinkurl{release_notes/v1.0-release-message.md}) are drafted;
(iii)~deposit on Zenodo to receive the DOI; (iv)~commit the
version-bump filling in the DOI in \nolinkurl{paper/main.tex},
\nolinkurl{CITATION.cff}, and the release-notes placeholders;
(v)~rebuild the PDF and snapshot to \nolinkurl{.stage/};
(vi)~tag and push; (vii)~create the GitHub Release.  This appendix
and the \nolinkurl{release_notes/README.md} parent describe the
full reproducibility chain from a fresh clone of any tagged
release.


%% ============================================================
\bibliographystyle{unsrt}
\bibliography{paper-bib/references}

\end{document}
